\begin{prob}
    In each of the problems below compute the average case time complexity (or
    expected time) of the given code. State your answer using asymptotic
    notation. Show your work for this problem by stating what the different
    cases are, the probability of each case, and how long each case takes. Also
    show the calculation of the expected time.

    \begin{subprobset}
        \begin{subprob}
            \begin{minted}[autogobble]{python}
                def foo(n):
                    # randomly choose a number between 0 and n-1 in constant time
                    k = np.random.randint(n)

                    if k > n/10:
                        for i in range(n):
                            print(i)
                    else:
                        print('Never mind...')

            \end{minted}

            \begin{soln}
                $\Theta(n)$.

                We can think of the two cases here as being 1) $k$ is randomly chosen
                to be greater than $n/10$, and 2) $k$ is randomly chosen to be less than $n/10$.
                The probability of the first case is $9/10$, and
                the probability of the second case is $1/10$.

                In the first case, linear time is taken since we must loop over \python{range(n)}.
                In the second case, constant time is taken since we print \python{"Never mind..."}
                and exit.

                Therefore, the average time taken is:

                \begin{align*}
                    \frac{9}{10} \Theta(n) + \frac{1}{10} \Theta(1)
                    =
                        \Theta(n) + \Theta(1)
                    =
                        \Theta(n)
                \end{align*}

            \end{soln}
        \end{subprob}


        \begin{subprob}
            \begin{minted}[autogobble]{python}
                def bogomax(numbers):
                    """Find the largest number by random guess and check."""
                    while True:
                        # randomly choose an element with uniform probability in constant time
                        guess = random.choice(numbers)

                        # check whether it is the largest
                        guess_is_a_maximum = True
                        for number in numbers:
                            if number > guess:
                                guess_is_a_maximum = False

                        if guess_is_a_maximum:
                            return guess
            \end{minted}

            In this part, you may assume that the numbers are distinct, and that $n$ is the size
            of \python{numbers}.

            Hint: if $0 < b < 1$, then $\sum_{p = 1}^\infty p \cdot b^{p-1} = \frac{1}{(1 - b)^2}$.

            \begin{soln}
                $\Theta(n^2)$.

                (Note in Fall 2022: originally the code in this problem had a typo where \python{guess_is_a_maximum = True}
                was \emph{inside} of the for-loop. This meant that, after the for-loop, \python{guess_is_a_maximum} did not
                actually tell us whether \python{guess} is a maximum; it only told us whether it was larger than the last
                number. This actually makes the analysis more difficult. The
                analysis below is for the corrected version of the problem.)

                Maybe the best way to divide this into cases is to consider Case 1 to
                be when we guess the right maximum on the very first try;
                Case 2 is when we guess the location of the max incorrectly on the
                first try, but get it right on the second; Case 3 is when we get the
                first two guesses wrong but get the third guess correct, and so on.
                In general, Case $\alpha$ is when we made incorrect guesses on the first
                $\alpha - 1$ iterations, but a correct guess on the $\alpha$th iteration,
                therefore returning.

                The probability of Case $\alpha$ is the probability of guessing wrong $\alpha - 1$
                times followed by guessing correctly. Since these guesses are all statistically
                independent, and since the probability of any one guess being wrong is
                $(n - 1)/n$ and the probability of it being right is $1/n$, we find that
                the probability of Case $\alpha$ is:
                \[
                    \underbrace{
                        \left( \frac{n-1}{n} \right)
                        \left( \frac{n-1}{n} \right)
                        \cdots
                        \left( \frac{n-1}{n} \right)
                    }_{\alpha - 1 \text{ times}}
                    \cdot \frac{1}{n}
                    =
                    \left(
                        \frac{n-1}{n}
                    \right)^{\alpha-1} \frac{1}{n}
                    =
                    \left(
                        1 - \frac{1}{n}
                    \right)^{\alpha-1} \frac{1}{n}
                \]

                The time taken to ``verify'' a guess (that is, the time taken for one iteration of the while-loop)
                is $\Theta(n)$, because we loop through all of the numbers and compare them to the guess.

                In case $\alpha$ we make $\alpha$ iterations, and so we take time $\alpha n$ in total.

                There are actually infinitely many cases here, since it is possible that
                we are consistently unlucky and never guess the right entry. Therefore,
                our summation will actually be an infinite series.

                Therefore the expected time over all possible cases is:
                \begin{align*}
                    \sum_{\alpha = 1}^\infty
                        P(\text{case } \alpha)
                        \cdot
                        T(\text{case } \alpha)
                        =
                    \sum_{\alpha = 1}^\infty
                    \left(
                        1 - \frac{1}{n}
                    \right)^{\alpha-1} \frac{1}{n} \cdot \alpha n
                    &=
                        \sum_{\alpha = 1}^\infty
                        \alpha
                        \left(
                            1 - \frac{1}{n}
                        \right)^{\alpha-1}
                        \\
                    \intertext{We can now use the hint. Letting $b = 1 - 1/n$, we have:}
                    &=
                        \sum_{\alpha = 1}^\infty
                        \alpha
                        b^{\alpha - 1}
                        \\
                    &=
                        \frac{1}{(1 - b)^2}
                    \intertext{Now substituting back in for $b = 1 - 1/n$:}
                    &=
                    \frac{1}{\left[1 - (1 - 1/n)\right]^2}
                        \\
                    &=
                        \frac{1}{1/n^2}
                        \\
                    &=
                        n^2
                \end{align*}

                So the average case time complexity is $\Theta(n^2)$.

            \end{soln}

        \end{subprob}

        \begin{subprob}
            \textbf{Extra Credit (3 points)}

            \textbf{Note}: this part is extra credit, and is totally optional! We
            suggest skipping it for now and coming back to it if you have time.

            The above version of \python{bogomax} can be improved by immediately
            returning when a maximum is found, as shown in the following code:

            \begin{minted}[autogobble]{python}
                import random

                def bogomax(numbers):
                    """Find the largest number by random guess and check."""
                    while True:
                        # randomly choose an element with uniform probability in constant time
                        guess = random.choice(numbers)

                        # check whether it is the largest
                        for number in numbers:
                            if number > guess:
                                break
                        else:
                            # yes, python has for-else blocks.
                            # if we get here, the for-loop completed without breaking
                            return guess
            \end{minted}

            This makes the code more performant, but harder to analyze.

            Analyze the average case time complexity of this code rigorously.
            You may assume that the elements of \python{numbers} are unique and
            that the maximum is equally-likely to be anywhere in the list.

            \begin{soln}
                $\Theta(n \log n)$.

                The difference between part b) and the part c) is that there is an if statement 
                which breaks out of the loop early. This means that every case is not always $\Theta(n)$.
                Remember that $\alpha$ represents the amount of iterations it takes until the max is guessed
                correctly. On the final iteration, the max goes through the entire loop without breaking so 
                it's $\Theta(n)$, but on the loops prior where there is a miss, we need to figure out the 
                expected time complexity. However, this is a bit difficult because we don't know which value 
                is selected or how the list looks.

                This will be the sub-problem that we need to solve for expected time complexity. We have a 
                case $\beta$ for each $n-1$ elements to find the expected time complexity. We define each case 
                $\beta$ as the $\beta$'th largest number selected. For every $\beta$'th case that means there 
                are $n - \beta$ larger numbers that if you run into will break you out of the loop. Finding the 
                probability of going through the list until running into a larger number is akin to drawing cards 
                until meeting an ace. The solution to that problem is described here:
                \url{https://aksoy.io/math/probability/book/2020/03/10/problem-40-first-ace.html}{here}. 

                That link describes the Principle of Symmetry, which states that $n$ randomly placed points
                will divide a segment into $n+1$ pieces, each of which has the same distribution. This means that
                for a selected random number, if we place it's larger numbers in the list, it will on average divide
                the list into segments of equal pieces. In the context of this problem, our number will on average 
                traverse the length of one segment before meeting a larger number and breaking out, which is 
                total length divided by \# of larger numbers $\frac{n}{n-\beta}$. 

                With this in mind, we can write down the expected number of iterations of the for-loop when
                our guess is the $\beta$th largest number:
                \begin{align*}
                    \sum_{\beta = 1}^{n-1}
                        P(\text{case } \beta)
                        \cdot
                        T(\text{case } \beta)
                    =
                    \sum_{\beta = 1}^{n-1}
                        \frac{1}{n-1} 
                        \cdot
                        \frac{n}{n-\beta}
                    =
                    \sum_{\beta = 1}^{n-1}
                        \frac{n}{(n-1)(n-\beta)}
                        \\
                    =
                    \sum_{\beta = 1}^{n-1}
                        \frac{n-1+1}{(n-1)(n-\beta)}
                    =
                    \sum_{\beta = 1}^{n-1}
                        \frac{n-1}{(n-1)(n-\beta)} 
                        + \frac{1}{(n-1)(n-\beta)} 
                        \\
                    = 
                    \sum_{\beta = 1}^{n-1}
                        \frac{1}{n-\beta} 
                        + \frac{1}{(n-1)(n-\beta)} 
                    \intertext{$\sum_{\beta = 1}^{n-1} \frac{1}{n-\beta}$ is  harmonic progression equivalent to 
                    $1 + \frac{1}{2} + \frac{1}{3} + ... + \frac{1}{n-1}$. There is no closed form solution, but when 
                    integrated is roughly $\log n$. The right hand side of the equation is equivalent except divided 
                    by $n-1$}
                    \approx
                    \log n + \frac{\log n}{n-1}
                    =
                    \Theta(\log n)
                \end{align*}

                Now that we know this, we know the time complexity of missing the max number grows at $\log(n)$. 
                However, this isn't the end. We still need to combine it with our part 2b). Instead of $\alpha n$ for 
                $T(\text{case } \alpha)$, we get $((\alpha - 1) \log(n) + n)$. This is because while the previous 
                $\alpha - 1$ cases have a time complexity of $\Theta(\log n)$ the last one will be the maximum 
                which has a time complexity of $\Theta(n)$

                \begin{align*}
                    \sum_{\alpha = 1}^\infty
                        P(\text{case } \alpha)
                        \cdot
                        T(\text{case } \alpha)
                    =
                        \sum_{\alpha = 1}^\infty
                        \left(
                            1 - \frac{1}{n}
                        \right)^{\alpha-1} \frac{1}{n} \cdot ((\alpha - 1) \log(n) + n)
                        \\
                    =
                        \sum_{\alpha = 1}^\infty
                        (\alpha - 1) \frac{\log n}{n} 
                        \left(
                            1 - \frac{1}{n}
                        \right)^{\alpha-1} +
                        \left(
                            1 - \frac{1}{n}
                        \right)^{\alpha-1}
                        \\
                    =
                        \frac{\log n}{n}
                        \sum_{\alpha = 1}^\infty
                            \alpha  
                            \left(
                                1 - \frac{1}{n}
                            \right)^{\alpha-1} -
                        \frac{\log n}{n}
                        \sum_{\alpha = 1}^\infty
                            \left(
                                1 - \frac{1}{n}
                            \right)^{\alpha-1}
                            +
                        \sum_{\alpha = 1}^\infty
                            \left(
                                1 - \frac{1}{n}
                            \right)^{\alpha-1}
                    \intertext{$\sum_{\alpha = 1}^\infty \left(1 - \frac{1}{n} \right)^{\alpha-1}$ will converge 
                    to some constant $c$, since it is a summation of geometric series with $0 < r < 1$. The left 
                    hand side can use the the same method already calculated in part $b)$ to evaluate to $n^2$}
                    =
                        \frac{\log n}{n} n^2 
                            -
                        \frac{\log n}{n} c
                            +
                        c
                    =
                        n \log n 
                            -
                        \frac{\log n}{n} c
                            +
                        c
                    = 
                        \Theta(n \log n 
                            -
                        \frac{\log n}{n} c
                            +
                        c)
                        = \Theta(n \log n)
                    \end{align*}
            \end{soln}

        \end{subprob}
    \end{subprobset}
\end{prob}
